\section*{Chapter \ref{ch:b2:retrosynthesis} --- Retrosynthesis and Multistep Strategy}

\begin{solution}{exo:b2:retrosynthesis:1}
\begin{itemize}
  \item 2-Phenylpropan-2-ol: propanone + \ce{PhMgBr}, or acetophenone + \ce{CH3MgBr}.
  \item Pentan-3-ol: propanal + \ce{CH3CH2MgBr}.
  \item 1-Cyclohexylethanol: ethanal + cyclohexylmagnesium bromide, or cyclohexanecarbaldehyde +
        \ce{CH3MgBr}.
  \item 2-Methylbutan-2-ol: propanone + \ce{CH3CH2MgBr}, or butanone + \ce{CH3MgBr}.
  \item 1-Phenylethanol: benzaldehyde + \ce{CH3MgBr}, or ethanal + \ce{PhMgBr}.
\end{itemize}
\end{solution}

\begin{solution}{exo:b2:retrosynthesis:2}
Acceptor \omterm{def:b2:retrosynthesis:disconnection}{synthon} $\mathrm{Ph\overset{+}{C}H{-}OH}$, equivalent benzaldehyde; donor \omterm{def:b2:retrosynthesis:disconnection}{synthon} \ce{CH3-},
equivalent methylmagnesium bromide (or methyllithium). (The other cut, with \ce{Ph-} as donor, uses
\ce{PhMgBr} and ethanal.)
\end{solution}

\begin{solution}{exo:b2:retrosynthesis:3}
$0.90 \times 0.85 \times 0.80 \times 0.95 \times 0.70 = 0.407$: about 41~\%.
\end{solution}

\begin{solution}{exo:b2:retrosynthesis:4}
\ce{LiAlH4}, needed to reduce the ester to the alcohol, would also reduce the ketone. Protect the
ketone as its cyclic acetal (ethane-1,2-diol, acid catalyst, water removed); reduce the ester with
\ce{LiAlH4}; remove the acetal with dilute aqueous acid: 5-hydroxypentan-2-one.
\end{solution}

\begin{solution}{exo:b2:retrosynthesis:5}
FGI: butane-1,3-diol $\Rightarrow$ 3-hydroxybutanal (reduction of the aldehyde, \ce{NaBH4}). Two-group
\omterm{def:b2:retrosynthesis:disconnection}{disconnection} of this $\beta$-hydroxy aldehyde at the $\alpha$--$\beta$ bond: two molecules of ethanal
(\omterm{def:b2:enolates-aldol:aldol}{aldol}). Forward: ethanal, dilute base, cold; then \ce{NaBH4}.
\end{solution}

\begin{solution}{exo:b2:retrosynthesis:6}
The cyclohexene ring is cut at the two \ce{C-C} bonds made in a \omterm{def:b2:frontier-orbitals:diels-alder}{Diels--Alder reaction}, those
allylic to the ring double bond on the side of the substituent: buta-1,3-diene and but-3-en-2-one, the
acetyl group on the \omterm{def:b2:frontier-orbitals:diels-alder}{dienophile}.
\end{solution}

\begin{solution}{exo:b2:retrosynthesis:7}
Numbering C1 (\ce{C=O}), C2=C3, then C4 with the two methyls
(\cref{met:b2:conjugate-additions:robinson-disconnection}): the \omterm{def:b2:enolates-aldol:aldol}{aldol} cut C2--C3 leaves an aldehyde
at C3 and a methyl ketone; the Michael cut C4--C5 leaves 2-methylpropanal (C3 its carbonyl, C4 its
$\alpha$ carbon) and but-3-en-2-one.
\end{solution}

\begin{solution}{exo:b2:retrosynthesis:8}
One step: benzene with butanoyl chloride and \ce{AlCl3} (\omterm{def:b2:aromatic-substitution:friedel-crafts}{Friedel--Crafts acylation}, no rearrangement,
single acylation). Two steps: benzaldehyde with propylmagnesium bromide, then oxidation of the secondary
alcohol. The Friedel--Crafts route is shorter and its reagents cheaper, but uses a full equivalent of
aluminium chloride, which ends in the aqueous waste.
\end{solution}

\begin{solution}{exo:b2:retrosynthesis:9}
Pyridinium chlorochromate in dichloromethane, the Swern oxidation or the Dess--Martin periodinane
give hexanal; acidified aqueous dichromate or permanganate give hexanoic acid, because in water the
aldehyde forms its hydrate, which is oxidised again (\cref{prop:b2:retrosynthesis:mild-oxidation}).
\end{solution}

\begin{solution}{exo:b2:retrosynthesis:10}
The amine reacts first, unprotected, if the alcohol is protected: protect the alcohol as a silyl ether
(it survives the acylation); acylate the amine; remove the silyl group with fluoride; acylate the
alcohol. If the amine had to be kept free while the alcohol reacts, it would be protected as a
\textit{tert}-butoxycarbonyl carbamate (removed by acid), orthogonal to the silyl ether (removed by
fluoride): either can be released without touching the other.
\end{solution}

\begin{solution}{exo:b2:retrosynthesis:11}
\ce{CH3COCH2CH2COCH3} is a 1,4-dicarbonyl: cut the central \ce{C-C} bond. The donor is the enolate of
ethyl 3-oxobutanoate (sodium ethoxide); the acceptor, an $\alpha$ carbon of a ketone, is provided by
1-bromopropan-2-one, \ce{BrCH2COCH3} (bimolecular substitution of bromide). Then aqueous hydrolysis of
the ester and heating remove \ce{CO2}: hexane-2,5-dione.
\end{solution}

\begin{solution}{exo:b2:retrosynthesis:12}
Retrosynthesis: the methyl ketone with a \ce{CH2} next to it is the product of the acetoacetic ester
synthesis; cut the bond between C3 and C4: the enolate of ethyl 3-oxobutanoate and the allylic halide
1-bromo-3-methylbut-2-ene, \ce{(CH3)2C=CHCH2Br}. Forward: sodium ethoxide in ethanol, then the halide;
then aqueous base, acidification and heating, which hydrolyse the ester and remove \ce{CO2}
(\cref{prop:b2:enolates-aldol:decarboxylation}): 6-methylhept-5-en-2-one.
\end{solution}

\begin{solution}{pb:b2:retrosynthesis:1}
\textbf{1.} \ce{HO-C6H4-CH2CH2COCH3}, the \ce{OH} para to the chain: a phenol and a ketone.
\textbf{2.} Hydrogenation of a \ce{C=C}: \ce{ArCH2CH2CO} $\Rightarrow$ \ce{ArCH=CHCO}, an
$\alpha,\beta$-unsaturated ketone.
\textbf{3.} (\textit{E})-4-(4-hydroxyphenyl)but-3-en-2-one, \ce{HOC6H4CH=CHCOCH3}.
\textbf{4.} An $\alpha,\beta$-unsaturated carbonyl: an \omterm{def:b2:enolates-aldol:aldol}{aldol condensation}, cut at the \ce{C=C}.
\textbf{5.} 4-Hydroxybenzaldehyde and propanone.
\textbf{6.} The aldehyde has no $\alpha$ hydrogen and can only be the electrophile; propanone, in
excess, is the only enolate source and reacts with the more reactive aldehyde rather than with itself.
\textbf{7.} (1) 4-hydroxybenzaldehyde, excess propanone, aqueous sodium hydroxide; acidification. (2)
\ce{H2} over palladium on carbon, at room temperature, until one equivalent is taken up.
\textbf{8.} The \ce{CHO} group stabilises the phenoxide, so the phenol is at least as acidic as phenol
($\mathrm pK_a$ 9.99). At $\mathrm{pH} > 13$, $[\ce{ArO-}]/[\ce{ArOH}] > 10^{13-10} = 10^3$: the
phenoxide.
\textbf{9.} The \ce{O-} pushes electron density through the ring onto the carbonyl carbon (a
resonance donor para to it): the aldehyde is less electrophilic, and the condensation slower.
\textbf{10.} A benzyl ether: the hydrogenation over palladium that reduces the \ce{C=C} also cleaves
the benzyl ether (hydrogenolysis), releasing the phenol in the same step.
\textbf{11.} Three: benzylation, \omterm{def:b2:enolates-aldol:aldol}{aldol condensation}, hydrogenation with deprotection.
\textbf{12.} No: under mild conditions (room temperature, palladium) the \omterm{def:b2:huckel:aromatic}{aromatic} ring is not reduced.
\textbf{13.} Not protecting: the phenoxide only slows the condensation, which a longer time or warming
compensate; one step saved is worth more than the gain in rate.
\textbf{14.} 4-Iodophenol and but-3-en-2-one.
\textbf{15.} \ce{HOC6H4I + CH2=CHCOCH3 + Et3N -> HOC6H4CH=CHCOCH3 + Et3NH+ + I-}, palladium catalyst.
\textbf{16.} At the terminal \ce{CH2}, the $\beta$ carbon: the aryl group inserts at the less hindered
end, and $\beta$-hydride elimination then gives the conjugated (\textit{E})-enone.
\textbf{17.} \omterm{def:b2:catalytic-cycles:oxidative-addition}{Oxidative addition} of the aryl iodide to \ce{Pd(0)}, coordination and insertion of the
alkene, $\beta$-hydride elimination, and regeneration of \ce{Pd(0)} by the base, which takes \ce{HI}
(\cref{ch:b2:catalytic-cycles}).
\textbf{18.} $162.188/(220.005 + 70.091 + 101.193) = 162.188/391.289 \approx 41.4~\%$.
\textbf{19.} \ce{C7H6O2 + C3H6O -> C10H10O2 + H2O}: $162.188/(122.123 + 58.080) = 162.188/180.203
\approx 90.0~\%$.
\textbf{20.} 100~\%: every atom of \ce{H2} and of the enone ends in the product.
\textbf{21.} $164.204/(122.123 + 58.080 + 2.016) = 164.204/182.219 \approx 90.1~\%$.
\textbf{22.} The \ce{C=C}: conjugated and unhindered, it binds and is hydrogenated on palladium far
faster than the ketone \ce{C=O}; stopping after one equivalent of hydrogen keeps the ketone.
\textbf{23.} \omterm{def:b2:enolates-aldol:aldol}{Aldol} route: cheap aldehyde, propanone and sodium hydroxide, water as by-product. Heck
route: an aryl iodide, a palladium catalyst, and but-3-en-2-one, very toxic.
\textbf{24.} The \omterm{def:b2:enolates-aldol:aldol}{aldol} route: two steps, cheap and safer reagents, atom economy above 90~\%.
\textbf{25.} It multiplies it by 0.90: $0.782 \times 0.90 \approx 70~\%$.
\textbf{26.} $0.85 \times 0.92 = \boldsymbol{\approx 78~\%}$ overall.
\end{solution}
